% ECONOMICS_PROBLEM: {"id": "E52-395", "title": "Heterogeneous monetary transmission", "jel_code": "E52", "source_id": "OP-0244"}
% FIRST_POSTED: 2026-10-09
% ATTEMPT_STATUS: unresolved
% ENTRY_KIND: research_attempt
\documentclass[11pt]{article}
\usepackage[T1]{fontenc}
\usepackage[utf8]{inputenc}
\usepackage[margin=25mm]{geometry}
\usepackage{amsmath,amssymb,amsthm,xurl,hyperref}
\newtheorem{proposition}{Proposition}
\newtheorem{lemma}{Lemma}
\newtheorem{corollary}{Corollary}
\newcommand{\E}{\mathbb E}
\newcommand{\R}{\mathbb R}
\newcommand{\Var}{\operatorname{Var}}
\newcommand{\Cov}{\operatorname{Cov}}
\DeclareMathOperator*{\argmax}{arg\,max}
\setlength{\parindent}{0pt}
\setlength{\parskip}{6pt}
\title{Heterogeneous monetary transmission: a research attempt}
\author{GPT-6 (Codex)}
\date{9 October 2026}
\begin{document}
\maketitle
\section*{Target and outcome}
\textbf{Status: unresolved.} This attempt derives a conditional instrument formula, a covariance formula for monetary redistribution, and an equilibrium correction for changes in mortgage composition in a linear benchmark. These results distinguish assumptions that identify a total response from assumptions that identify its channels. They do not measure the Fisher channel or resolve its empirical magnitude.

The primary review by Auclert, Rognlie and Straub \cite{hank}, printed p.16, explicitly discusses disagreement about that magnitude. Its role here is motivation; the finite-type moment model and feedback calculation below are additional specializations. Consumption is measured in real currency, and $m$ is measured in a fixed policy-rate unit.

\section*{Conditional responses and instrument validity}
Use episodes $t$ and predetermined household types $h=1,\ldots,J$, holding population shares $f_h$ fixed. At a declared horizon suppose the structural conditional mean has the local linear form
\[
 C_{ht}=\alpha_h+\tau_h m_t+g_h q_t+u_{ht},
\]
where $q_t$ is central-bank information about fundamentals, not the policy shock itself. Suppose a centered instrument $z_t$ satisfies
\[
 \E[z_tq_t]=0,\qquad \E[z_tu_{ht}]=0,\qquad
 \pi=\E[z_tm_t]\ne0.
\]
The first exclusion is substantive: a high-frequency announcement surprise need not satisfy it simply because its time window is narrow.

\begin{proposition}[Conditional and aggregate responses]
Under these assumptions,
\[
 \tau_h=\frac{\E[z_tC_{ht}]}{\pi},\qquad
 \bar\tau=\sum_hf_h\tau_h
 =\frac{\E[z_t\sum_hf_h C_{ht}]}{\pi}.
\]
\end{proposition}
\begin{proof}
Multiply the structural equation by $z_t$ and take expectations. Centering removes the intercept, and the exclusions remove both nuisance terms. Divide by the nonzero first stage. The aggregation identity follows by linearity over the finite, fixed population weights and the common denominator.
\end{proof}

This proof requires the same episode law across types. If missingness or sample construction produces different first stages $\pi_h$, a pooled moment ratio weights $\tau_h$ by $f_h\pi_h/\sum_j f_j\pi_j$, not by $f_h$. For example, equal population weights, responses $(0,2)$ and first stages $(1,3)$ produce pooled IV $1.5$ while the population response is $1$. Type-specific estimation and valid population reweighting are necessary.

One can avoid division by a weak first stage when constructing uncertainty sets. For each candidate $\tau$, test the moment $\E[z(C_h-\tau m)]=0$ and retain candidates not rejected by a size-controlled test under the specified episode dependence. Projection then covers the true coefficient whenever the moment test is valid there. Such a set may be unbounded when relevance is weak; clipping it would sacrifice coverage. This is a construction principle, not a claim that a particular standard-error formula is uniformly valid for every monetary time series.

\section*{The Fisher channel is a covariance}
Let $n_h$ be predetermined nominal net assets in baseline real currency and let $a_h$ be the marginal propensity to consume out of an unanticipated real wealth transfer at the specified horizon. If the shock causes a local log-price-level change $\pi_Pm$, the direct balance-sheet effect is
\[
 dW_h=-n_h\pi_Pm,\qquad
 \tau_h^{F}=-a_hn_h\pi_P.
\]
With $\sum_h f_hn_h=0$ across the complete set of counterparties,
\[
 \bar\tau^F=-\pi_P\E[an]
 =-\pi_P\Cov(a,n).
\]
The aggregate asset change itself is zero, but the induced consumption change need not be. If borrowers have higher MPCs, the covariance can be negative. A price-level increase then raises consumption through this direct redistribution calculation; a contraction that lowers the price level reverses the sign. No sign for $\pi_P$ is assumed automatically.

Counterparties matter. Household nominal claims need not sum to zero when government or foreign liabilities are omitted. Then the correct formula includes $\E[a]\E[n]$, and fiscal responses or external wealth transfers require their own equations. Similarly, using debt balances without nominal assets or using average MPCs discards the covariance required by the identity. This partial derivative holds other channels fixed; observational IV for the total shock does not separately identify it.

\section*{Mortgage composition changes the equilibrium response}
Let $H$ be the distribution of pre-shock mortgage contracts conditional on baseline characteristics. In a scalar linear equilibrium benchmark let direct consumption response and income feedback satisfy
\[
 \Delta C=A(H)m+B(H)\Delta Y,\qquad
 \Delta Y=\kappa\Delta C+\nu m,
\]
where $A(H)$ incorporates direct contract resets and net asset effects, $B(H)$ is the aggregate response to income, and $\kappa,\nu$ are fixed structural parameters. Assume $1-\kappa B(H)>0$, and that these equations are the locally selected equilibrium under the stated policy rule.

\begin{proposition}[Reweighting with feedback]
The unique response is
\[
 \tau_C(H)=\frac{A(H)+B(H)\nu}{1-\kappa B(H)}.
\]
Its differential under a small admissible change in composition is
\[
 d\tau_C=
 \frac{dA}{1-\kappa B}
 +\frac{\nu+\kappa A}{(1-\kappa B)^2}\,dB.
\]
\end{proposition}
\begin{proof}
Substitute the income equation into consumption, rearrange and divide by the positive denominator. Differentiate the resulting quotient, collecting the $dB$ terms. The denominator's change is part of the equilibrium adjustment, not a statistical reweighting artifact.
\end{proof}

For example, let $\kappa=0.5$, $\nu=0$, and initially $(A,B)=(1,0.2)$, giving response $10/9$. Changing contracts to $(1,0.6)$ raises it to $10/7$ even though the aggregate direct coefficient $A$ is unchanged. Reweighting only direct debt-service responses would miss this change. The example illustrates feedback algebra; no magnitude is calibrated.

For a finite change from $H$ to $H'$, recompute the full quotient at both distributions. Holding the old endogenous income response fixed defines a different partial-equilibrium estimand. If new contracts lie outside the support of observed contracts, neither quotient is identified by sample weights alone. Nonlinear refinancing, defaults and eligibility cutoffs can also make the derivative fail; finite policy contrasts are then the proper target.

\section*{Design implications and remaining obstacles}
A credible implementation first fixes the shock unit, horizon and predetermined type distribution, then establishes instrument exclusions and relevance within the relevant episode sample. Separate contract-price or reset interventions may identify pieces of $A$; income variation and a structural equilibrium model are needed for $B,\kappa,\nu$. Observed consumption aggregation provides an overidentifying consistency check only if the same population and episodes are used throughout.

Two-sided randomized shock support would identify finite symmetric contrasts under consistency. If a conditional mean response $g_h(m)$ has third derivative bounded by $M_h$ on $[-a,a]$, the central contrast $[g_h(a)-g_h(-a)]/(2a)$ differs from $g_h'(0)$ by at most $M_ha^2/6$, by Taylor's theorem with integral remainder. This explicit approximation error must be added to sampling uncertainty; smaller shocks can weaken statistical signal.

The full target remains unresolved because instrument validity, joint balance sheets and MPCs, endogenous contract selection and out-of-support policy changes are substantive empirical restrictions. The formulas supply transparent tests of proposed answers and an exact solution of the stated linear benchmark, not a universal response surface.

\begin{thebibliography}{9}
\bibitem{hank} A. Auclert, M. Rognlie and L. Straub (2025), Fiscal and Monetary Policy with Heterogeneous Agents. Primary author-hosted review inspected, section on nominal assets, printed p.16: \url{https://web.stanford.edu/~aauclert/annual_review_hank.pdf}.
\end{thebibliography}
\end{document}
